\documentclass[11pt]{article}
\usepackage{amsmath,amssymb,amsthm,fullpage,hyperref}
\newtheorem{proposition}{Proposition}
\title{Positive Mixtures Do Not Remove the Signed-Span Problem}
\author{}
\date{November 24, 2026}
\begin{document}
\maketitle

\paragraph{Status.}
This is a written audit of a possible entropy-based route to the
matching lower bound. It does not establish that lower bound or an
upper-bound improvement. The elementary reduction below applies to
all affine numerators and all everywhere-positive affine denominators,
so in particular to every admissible head. No model definition is
changed. The reduction is not Lean checked.

\section{The literature theorem and its scope}
Gassiat and van Handel, \emph{The local geometry of finite mixtures},
\url{https://arxiv.org/abs/1202.3482}, Theorems 3.1 and 3.3,
give local bracketing entropy estimates for convex mixtures of
translates of a smooth density. Their exponent is linear in the
number of components. This is relevant because a direct parameter
count for rational sums loses a logarithm of the number of heads.
However, the constants depend on the reference mixture; page 10
explicitly leaves uniform control over that reference unresolved.
Moreover, the local geometry theorem uses the identifiable structure
of location mixtures. Remark 3.6 identifies that theorem as the part
requiring a new argument for other dictionaries. Thus this result
does not supply a uniform signed-mixture or strict-sign bound for
HeadComplexity.

\section{Every signed score is already a positive-mixture direction}
Let $X$ be a nonempty finite set. Suppose $B_h(x)>0$ on $X$ and
write
\[
 F(x)=c+\sum_{h=1}^H\frac{A_h(x)}{B_h(x)},\qquad H\geq1.
\]
Absorb the bias by setting $C_h=A_h+(c/H)B_h$, so
$F=\sum_h C_h/B_h$. All $C_h$ remain affine when $A_h,B_h$
are affine.

\begin{proposition}
There is $\epsilon_0>0$ such that, for every
$0<\epsilon<\epsilon_0$, the functions
\[
 g_{h,\epsilon}(x)=
 \frac{B_h(x)+\epsilon H C_h(x)}{B_h(x)}
\]
are strictly positive on $X$, and their equal-weight mixture obeys
\[
 \frac1H\sum_h g_{h,\epsilon}=1+\epsilon F.
 \tag{1}\label{eq:exact}
\]
The denominators, including all their slope signs, are unchanged.
For any norm on $\mathbb R^X$, if $F\ne0$, the normalized direction
from the constant function is exactly
\[
 \frac{H^{-1}\sum_h g_{h,\epsilon}-1}
 {\|H^{-1}\sum_h g_{h,\epsilon}-1\|}
 =\frac F{\|F\|}.
\]
\end{proposition}
\begin{proof}
Let $M=\max_{h,x}|HC_h(x)/B_h(x)|$, which exists by finiteness.
Take $\epsilon_0=1/(1+M)$. Then
$g_{h,\epsilon}(x)=1+\epsilon HC_h(x)/B_h(x)>0$.
Summing gives \eqref{eq:exact}; the normalized identity follows
because $\epsilon>0$.
\end{proof}

Consequently replacing signed affine fractions by positive affine
fractions near a constant does not simplify the required family of
normalized directions. The entire signed score class occurs there
already. In particular, numerator positivity cannot justify a
second-order-only cancellation argument: the reference fraction
$B/B=1$ is available for every denominator $B$.

If probability densities are required, let $\mu$ be the uniform
probability measure on $X$, and write $\bar F=\mathbb E_\mu F$.
Normalize the mixture in \eqref{eq:exact} by its positive mean.
Its density $f_\epsilon$ satisfies
\[
 f_\epsilon-1
 =\frac{\epsilon}{1+\epsilon\bar F}(F-\bar F).
\]
Thus every centered signed score occurs as an exact normalized
direction of positive affine-fraction mixture densities. Individual
component normalization merely changes their positive mixing weights.
The unrestricted threshold restores the discarded scalar mean.

\section{What a usable new theorem would have to control}
On $X=\{0,1\}^n$, after denominators are fixed, every score lies
in the linear space generated by $1$ and
\[
 \frac1{B_h(x)},\qquad \frac{x_i}{B_h(x)}
 \quad (1\leq h\leq H,\ 1\leq i\leq n).
\]
Its dimension is at most $1+(n+1)H$. The open problem is the union
of strict sign patterns of these spaces as all admissible
denominators vary. Positivity of the fractions near $1$ leaves this
union intact by the proposition. A uniform entropy estimate for
that union would be substantive new input, rather than a direct
application of the location-mixture theorem. Metric entropy alone
would also require a separate argument converting normalized
directions into strict sign counts without a prescribed margin.

\end{document}
